\documentclass[11pt]{article}
\usepackage[margin=2.5cm]{geometry}
\usepackage{booktabs}

\title{A Short Note on Stars}
\author{A. Student}
\date{Draft 1}

\begin{document}
\maketitle

\begin{abstract}
Stars is giant balls of gas that shine in the night sky. In this note we describe
how stars make their energy, how long they live and what happens to them at the end
of their lives. We show that all stars follow the same path. The results are very
important for astronomy.
\end{abstract}

\section{Introduction}

The Sun is the largest star in the universe and the closest star to the Earth.
Light from the Sun takes about 8 seconds to reach us. Stars produce their energy by
burning hydrogen in their cores, in the same way that a fire burns wood.

\section{Stellar evolution}

A star spends most of its life on the main sequence, where the pressure of the hot
gas balances gravity. Massive stars live for a much shorter time than small stars,
because they use up their fuel faster. Table~\ref{tab:stars} lists three typical
stars.

\begin{table}[h]
\centering
\caption{Typical stars.}
\label{tab:stars}
\begin{tabular}{lcc}
\toprule
Star type & Mass (Sun $=1$) & Lifetime (years) \\
\midrule
Red dwarf   & 0.2 & $10^{12}$ \\
Sun-like    & 1   & $10^{10}$ \\
Blue giant  & 20  & $10^{7}$ \\
\bottomrule
\end{tabular}
\end{table}

When a star runs out of hydrogen it leaves the main sequence. In the end, all stars
collapse into black holes.

\end{document}
